\section{Introducción}\label{sec:intro}

FIRA financia y asegura siembras de cebada de temporal, la que depende solo de la lluvia. Hidalgo, Tlaxcala y Puebla
están entre los cinco estados con mayor producción nacional de cebada grano~\cite{sader2025}. FIRA necesita saber qué
rendimiento esperar, y con qué riesgo, en lugares donde no hay parcelas medidas.

Las etiquetas del reto no sirven para eso. Se dividen por municipio en dos grupos que promedian 4.49 y 3.05\,t/ha.
Los rendimientos de un grupo se registran con dos decimales y los del otro están redondeados. Un reporte técnico
interno del equipo~\cite{regimenes} probó clima, altitud, suelo, fenología por satélite y cercanía al Popocatépetl,
y ninguna de esas variables explica la diferencia de 1.44\,t/ha entre los grupos. Por eso usamos el rendimiento
oficial del Servicio de Información Agroalimentaria y Pesquera (SIAP)~\cite{siap}. El SIAP publica una cifra por
municipio con el mismo procedimiento de reporte. Eso no asegura que se haya medido igual en campo, y el SIAP a veces
copia una misma cifra entre municipios (Sección~\ref{sec:datos}). El modelo predice el rendimiento promedio de un
municipio según esa cifra oficial, no el de una parcela.

\textbf{Resultados principales.}
\begin{itemize}
\item \textbf{La cebada rinde más donde el verano es fresco.} Cada grado más de temperatura máxima en julio y agosto
se asocia con cerca de 0.8\,t/ha menos. Es una asociación y no prueba que el calor sea la causa.
\item \textbf{Ese dato basta para estimar lugares sin datos.} Lejos de municipios con SIAP, el modelo de temperatura
se equivoca menos que copiar a los vecinos, 0.44 contra 0.57\,t/ha. Como elegimos la variable viendo resultados, una
selección repetida dentro de la validación reduce la ventaja de 0.13 a cerca de 0.07\,t/ha.
\item \textbf{El suelo y el clima del año no ayudan.} El suelo no mejora el modelo y el clima mensual del año no
anticipa la cosecha.
\end{itemize}
